\subsection{Eckhoff \& Behler (2021): spin-dependent HDNNP for MnO}
\label{sec:appendix-example-eckhoff2021spin}

\paragraph{Q1 -- Bibliographic identification.}
Eckhoff and Behler extend the high-dimensional neural network potential
framework with spin-dependent atom-centered symmetry functions (sACSFs)
to capture magnetic systems. Published in \emph{npj Comput.\ Mater.} 7,
170 (2021); arXiv:2104.14439.
\lstinputlisting[language=turtle]{artifacts/kg/listings/eckhoff2021spin/q01-bibliographic.ttl}

\paragraph{Q2 -- Material system.}
The headline material is rock-salt MnO (and a vacancy-rich
Mn$_{0.969}$O variant). The AFM-II rhombohedrally distorted ground
state with N\'eel temperature ${\sim}116$\,K is the focus.
\lstinputlisting[language=turtle]{artifacts/kg/listings/eckhoff2021spin/q02-material.ttl}

\paragraph{Q3 -- Reference calculation method.}
Reference data come from spin-polarised DFT calculations.
\lstinputlisting[language=turtle]{artifacts/kg/listings/eckhoff2021spin/q03-reference-method-type.ttl}

\paragraph{Q4 -- Reference settings.}
Calculations use FHI-aims (version 200112.2) with the HSE06
screened-hybrid functional ($\omega = 0.11\,a_0^{-1}$), all-electron
numerical atom-centred orbitals (intermediate basis, excluding
auxiliary 5g hydrogenic functions), and a $\Gamma$-centred
$2{\times}2{\times}2$ k-mesh on $2{\times}2{\times}2$ MnO supercells.
\lstinputlisting[language=turtle]{artifacts/kg/listings/eckhoff2021spin/q04-reference-settings.ttl}

\paragraph{Q5 -- Training dataset.}
The reference set contains 3{,}101 supercells across magnetic states:
1{,}387 MnO + 1{,}421 Mn$_{0.969}$O for training, 156 MnO + 137
Mn$_{0.969}$O for testing. Energies and forces are covered.
Provenance is \emph{in-house} (assembled for this work).
\lstinputlisting[language=turtle]{artifacts/kg/listings/eckhoff2021spin/q05-dataset.ttl}

\paragraph{Q6 -- Sampling strategies.}
Four strategies were combined: magnetic-state enumeration (FM, AFM-I,
AFM-II, plus excited spin orderings); atomic-displacement
perturbations from ideal lattice positions; lattice-parameter
perturbations to capture the rhombohedral distortion; and Mn-vacancy
sampling for off-stoichiometric configurations.
\lstinputlisting[language=turtle]{artifacts/kg/listings/eckhoff2021spin/q06-sampling.ttl}

\paragraph{Q7 -- MLIP method, components, implementation.}
The method (mHDNNP) is a Behler--Parrinello HDNNP with three hidden
layers (20/15/10 neurons per element) consuming spin-augmented ACSFs:
radial $M^0/M^+/M^-$ and angular
$M^{00}/M^{++}/M^{--}/M^{+-}$ functions. Implementation is RuNNer
v1.00 (modified to support sACSFs); n2p2 is used for LAMMPS MD.
\lstinputlisting[language=turtle]{artifacts/kg/listings/eckhoff2021spin/q07-method.ttl}

\paragraph{Q8 -- Hyperparameter settings.}
Reported settings: cutoff radius $r_c = 10.5\,a_0$ (Bohr); ACSF
amplitude threshold $M_{\text{thresh}} = 0.25$; per-element
architecture three layers $\times$ (20, 15, 10) neurons.
\lstinputlisting[language=turtle]{artifacts/kg/listings/eckhoff2021spin/q08-hyperparameter-settings.ttl}

\paragraph{Q9 -- MLIP run and trained model.}
A single training run produces the production mHDNNP for MnO.
\lstinputlisting[language=turtle]{artifacts/kg/listings/eckhoff2021spin/q09-run-and-model.ttl}

\paragraph{Q10 -- Benchmark results.}
Test-set RMSEs: 1.11\,meV/atom on cohesive energies and
0.066\,eV/\AA{} on atomic-force components. The paper also reports
phonon and N\'eel-temperature predictions but in derived form rather
than as separate $\BenchmarkResultC$ instances.
\lstinputlisting[language=turtle]{artifacts/kg/listings/eckhoff2021spin/q10-benchmark-results.ttl}

\paragraph{Q11 -- Computational resources.}
\emph{Not reported.}
\lstinputlisting[language=turtle]{artifacts/kg/listings/eckhoff2021spin/q11-resources.ttl}

\paragraph{Q12 -- Gaps.}
Beyond Q11, the most consequential gap is the lack of first-class
ontology coverage for \emph{spin-dependent} descriptors and magnetic
configurations. The mHDNNP's spin-augmentation is currently captured
only through prose in the $\dlConcept{FunctionalForm}$ label and the
magnetic-state-sampling strategy; the ontology has no concept for
collinear/non-collinear spin states, magnetic-moment-per-atom
properties, or spin descriptors. Other unreported items: the FHI-aims
plane-wave-equivalent energy cutoff (replaced by basis-set choice in
NAO); explicit $\dlRole{frozenCore}$; and any HSE06 numerical
parameters beyond the screening parameter $\omega$.
